% 05-ablation-analysis.tex — truth: 本文件（2026-09-10 交割：R22 后全部正文修订直改 tex，原装配源 md 已冻结为历史设计稿，见其头部治理注记）
% 纪律: E5 探索性发现措辞按裁决①(效率主张禁用); tab3 表格 + 低资源梯度段(成对呈现纪律: 低资源档与饱和档必须成对出现);
% fig4 caption 逐字取 FIGURE_SOURCE.md; 四步分析结构 observe→interpret→implicate→next。
\section{Ablation and analysis}\label{sec:ablation}

Table~\ref{tab:ablation} summarizes the component-level ablations; each row removes, replaces, or perturbs one design element and reports the measured consequence under its disclosed caliber.

\begin{table}[pos=htbp]
\centering
\caption{Component ablations. Every number is tier-tagged; tiers are never mixed. Paired designs are described inline. Device calibers are disclosed where a row spans devices (row 1: gradient tiers CPU, saturation tier GPU); cross-device absolute accuracies are not poolable (Section~\ref{sec:conclusion}, L10).}
\label{tab:ablation}
\footnotesize
\hyphenpenalty=10000
\begin{tabularx}{\linewidth}{>{\raggedright\arraybackslash}p{2.6cm}>{\raggedright\arraybackslash}p{1.8cm}>{\raggedright\arraybackslash}X}
\hline
Switch & Claim & Result (caliber) \\
\hline
$-$ self-supervised pretraining & C3 & Benefit decays to statistical equivalence with budget: $+21.2$\,pp at 88 clips (scratch $\approx$2.3$\times$ random), $+9.9$\,pp at 176, $-2.3$\,pp at 352 (crossover), $+0.15$\,pp n.s.\ at 1760 \emph{(synthetic; scratch vs.\ warm, 3 seeds; gradient tiers CPU, saturation GPU)} \\
$-$ semantic warm-start (generic-SSL instead) & C6 & Corrected end-to-end protocol (final head on seed truth + pool pseudo-labels; 10 seeds at 2 and 4 clips/class; full-budget arm deterministic): both arms near chance at 2 clips/class (top-1 9.8 vs.\ 15.2; macro-F1 5.3 vs.\ 10.4, Holm-corrected $p=0.023$ favoring generic-SSL), same ordering at 4 (top-1 10.9 vs.\ 15.9, n.s.; macro-F1 7.3 vs.\ 13.9, corrected $p=0.049$); at full supervision the reference still favors the task-pretrained backbone (33.93 vs.\ 30.36). The earlier warm-start advantage at low budget was an artifact of a true-label head and is superseded \emph{(public-real; AimCLR@InterPet4D, mocap domain gap disclosed)} \\
$-$ unsupervised segmentation (replace SMQ) & C4 (auxiliary) & SMQ 0.458$\pm$ 0.049 $>$ uniform window 0.399$\pm$ 0.035 $>$ random-cut null 0.323$\pm$ 0.022; all three pre-registered IoU gates pass. A fixed-grid probe (0.453$\pm$ 0.027) is on par with SMQ, and uniform windows lead on boundary F1 (0.396 vs.\ 0.343)---no consistent SMQ advantage. SMQ is positioned as an auxiliary proposal source; a downstream ablation swapping SMQ's proposals for equal-count fixed windows under the same frozen-feature head measures identical clip-level accuracy (45.95\% vs.\ 45.95\%, $\Delta = 0.0$) \emph{(public-real; seeds pseudo-GT, 4 episodes)} \\
$-$ anchor guidance & C5 & Prototype purity drops from 0.5339 to a 0.3306 random-assignment baseline ($1.61\times$; all-majority baseline 0.4858); seed-ratio 25\%$\to$100\% adds $\leq$1\,pp at $K\in\{5,10\}$ but costs 2.4\,pp at the best sweep cell $K{=}14$ \emph{(public-real / P0.3)} \\
$-$ pseudo-label iteration (single round) & C5 & Pool precision 0.5125 $\to$ 0.6913 ($+17.9$\,pp at round 1, a point estimate); the pre-registered first-to-last comparison over all iterated rounds is direction-consistent but \textbf{not significant after Holm--Bonferroni correction} (corrected $p=0.090$ over the three-test family; raw $p=0.030$); the precision-drop stopping rule fired after two degraded rounds, selecting round~1. We claim \emph{no robust benefit from multiple iterations}: round-1 is the reported operating point, multi-round gains directional only, a dedicated multi-round ablation registered as future work \emph{(public-real / P0.4, physics-prior consensus caliber)} \\
$-$ uncertainty sampling (use random) & C7 (exploratory) & No entropy advantage under cold-start weak scorers: random leads in mean by $+7.9$\,pp at budget 100 (2/3 seeds) and $+7.1$\,pp at 200 (3/3 seeds); direction confirmed with strong in-domain scorers ($+4.2$--$5.0$\,pp at every budget above the shared initial set) \emph{(synthetic; negative result)} \\
seed-noise injection \{10/20/30\%\} & C5 & Purity 0.5339 $\to$ 0.5267 $\to$ 0.5201 $\to$ 0.5025 at 30\% ($-3.1$\,pp, still $1.52\times$ baseline)---iteration does not amplify seed noise at this scale \emph{(public-real / P0.3)} \\
\hline
\end{tabularx}
\end{table}

\paragraph{Pretraining value is a function of annotation scarcity.}
Across four matched budgets (88/176/352/1760 clips, scratch vs.\ pretraining-initialized, three seeds each), the fine-tuning benefit decays to statistical equivalence: $+21.2$\,pp at 88 clips (scratch $\approx$2.3$\times$ random, warm $7\times$), $+9.9$\,pp at 176 (warm-start std $3.8\times$ narrower), $-2.3$\,pp at 352 (crossover region, within one arm's std), $+0.15$\,pp (n.s.) at 1760. The zero crossing lies between 10 and 20 clips per class; the bracket is the claim, and low-resource and saturated tiers are reported paired throughout. Together with the $k$-NN probe of Section~\ref{sec:main-results} (E1), this gives two complementary lines of evidence for C3: representation discriminability (probe) and fine-tuning value (gradient). Mechanism: the scarcer the labels, the less a discriminative representation can be learned from data, and the cross-domain prior of the pretraining-initialized encoder becomes the dominant usable structure at the smallest budgets.

\paragraph{Segmentation component status (downgraded).}
SMQ's IoU advantage over the uniform-window arm (0.458 vs.\ 0.399) does not extend to boundary F1 (0.343 vs.\ 0.396), and the fixed-grid probe is on par (0.453 $\pm$ 0.027); evaluation uses seed-derived pseudo-ground truth over four episodes. A dedicated downstream recognition ablation isolates SMQ's contribution directly: on the public-real tier (7-class seed-pseudo-GT caliber, 151/74 train/val clips), replacing SMQ's proposals with equal-count fixed windows under the identical frozen-feature linear head yields identical clip-level accuracy (45.95\% vs.\ 45.95\%, $\Delta = 0.0$). SMQ is therefore treated throughout as an auxiliary proposal-generation step inherited from prior work~\citep{gokay2025smq}, not a validated contribution; the claims of Section~\ref{sec:main-results} do not rest on it.

\paragraph{Multi-round pseudo-label iteration: no robust benefit.}
The round-1 endpoint gain in pool precision ($0.5125 \to 0.691$, $+17.9$\,pp) is real but is a point estimate; the pre-registered paired first-to-last comparison over all iterated rounds yields $\Delta$ $+10.69 \pm 3.28$\,pp with raw $p=0.030$, \emph{not} significant after Holm--Bonferroni correction over the three-test family (corrected $p=0.090$). The precision-drop stopping rule fired after two consecutive degraded rounds, selecting round~1 as the reported operating point. We make no claim that multiple iterations improve the pipeline: all reported operating points use round~1, the first-to-last trend is direction-consistent only, and a controlled multi-round ablation is registered as future work.

\paragraph{Semantic warm-start versus generic-SSL on real skeletons.}
The synthetic-tier row above asks whether pretraining helps at all; on the public-real tier we ask the sharper question of \emph{which} pretraining earns the low-label budget. Holding the entire end-to-end pipeline fixed (seeds, anchor-guided pseudo-labeling, head, splits, leakage guards) and replacing the backbone initialization together with its bound native input preprocessing (each arm uses the preprocessing its method was trained with; disclosed)---task-pretrained semantic warm start versus a generic contrastive encoder pretrained on motion-capture dog skeletons (AimCLR)---the corrected pseudo-label protocol finds both arms at or near chance at 2 clips/class (top-1 9.8\% vs.\ 15.2\%; chance 11.1\%), with the generic-SSL arm ahead on macro-F1 (10.4\% vs.\ 5.3\%, per-tier Holm-corrected $p=0.023$) and the same ordering at 4 clips/class; at full supervision the deterministic supervised reference still favors the task-pretrained backbone (33.93\% vs.\ 30.36\%; Section~\ref{sec:main-results}, E7c). The generic arm carries a mocap-to-in-the-wild domain gap its original evaluation did not: it is a transfer lower bound, not a best-tuned AimCLR. The reading is a boundary condition rather than a mechanism: at a 13\% budget on this tier, neither initialization yields pseudo-labels good enough to train on, and the earlier warm-start advantage---obtained under a protocol that trained the final head on pool ground truth---does not survive the correction. What the tier's supervised comparisons support (E4 pool precision, direction-consistent but corrected $p=0.090$; the five-arm study in Section~\ref{sec:comparison}) is that task-pretrained features are better when labels are present. The released per-arm pool-precision diagnostics ($0.11$ warm vs.\ $0.15$ generic-SSL) suggest pseudo-label quality does \emph{not} track supervised feature quality at budget---a post-hoc observation, not a controlled test.

\paragraph{Canine rescue levers: three pre-registered nulls.}
The alignment, backbone-adaptation, and budget levers below were pre-registered to rescue the canine tier's low-budget failure (Section~\ref{sec:main-results}, E7). All three returned null, and we consolidate them here as the boundary-condition evidence referenced in Section~\ref{ssec:boundary}.

\paragraph{Label-alignment force: a pre-registered null on the canine tier.}
The E7/E7b negative boundary raises the question of whether a stronger pseudo-label correction can rescue the canine tier at a 13\% budget. We pre-registered (PSD-ALIGN-PREREG-001, frozen before the run) three ground-truth-free alignment mechanisms against the R16 control, all sharing the corrected protocol and differing only in how the seed-anchored prototype geometry constrains pseudo-labels: V1 requires head--prototype top-1 agreement for pool entry, V2 adds a seed-prior per-class quota, V3 lets the prototype path assign throughout. Over ten seeds on AK v1 spc2, the best variant (V2) reaches 14.1\% $\pm$ 6.7 top-1---winning 8 of 10 same-seed pairs against the control (9.8\%, mean $+4.3$\,pp, Wilcoxon raw $p=0.13$, not significant)---and V3 roughly doubles pool pseudo-label precision (0.21 vs.\ control 0.11)---but the pre-registered rescue bar (mean $\geq 20\%$) is not met, so the decision rule returns \textbf{NULL}. We report this as a bounded negative (the first of three null rescue levers, Section~\ref{ssec:boundary}): the alignment force improves pseudo-label quality and is directionally consistent, but on this tier the binding constraint is the seed-anchored feature geometry itself (r0 prototype precision $\approx$0.30 caps what any gate can select), not the gating rule---consistent with the label-in-the-loop diagnosis of Section~\ref{sec:conclusion} (L9). Evidence: the released experiment ledger.

\paragraph{Two further levers: backbone adaptation and label budget, both null.}
If gating was not the binding constraint, two orthogonal levers remain: a stronger backbone and more labels. We pre-registered both (PSD-DAP-PREREG-001, PSD-BUDGET-PREREG-001; ADR 0008). \emph{Backbone}: domain-adaptive continued pretraining (two-view InfoNCE) on 503 unlabeled APTv2 canid skeleton clips, initialized from the physics-layer backbone. A pre-registered fidelity self-check fired first---the full-supervision reference on AK dropped from 33.93\% to 25.0\% ($<$30\% line), i.e.\ catastrophic forgetting---and the adapted backbone did not help downstream (V2+DAP 11.4\% vs.\ 14.1\% for the P1 gate, paired wins 3/10), corroborating that APTv2$\to$AK is bottlenecked by topology/appearance mismatch rather than adaptation strength. \emph{Budget}: sweeping the seed budget from spc2 (13\% labels) to spc12 (77\% labels), the best variant (V2) rises only from 14.1\% to 22.1\% and never reaches the 24\% (70\%-retention) knee at any budget; the pseudo-label pool (precision 0.04--0.17) adds no measurable value over the seed-only arm. Three independent levers---gating, backbone, budget---thus agree: on the AK self-extraction tier the end-to-end low-resource claim is not rescued by any lever we tested, and the binding constraint is the tier's data quality itself; the three nulls jointly constitute the boundary condition analyzed in Section~\ref{ssec:boundary}. The pre-registered real-K9 pilot (Section~\ref{sec:datasets}) is designed to test this falsification on target-domain data of adequate quality. Evidence: the released experiment ledger.

\paragraph{Uncertainty sampling: an exploratory negative finding (C7, downgraded by pre-registration).}
As an exploratory finding, uncertainty sampling did not help in our regime. Under cold-start weak scorers on the synthetic tier, entropy sampling showed no detectable advantage: random led in mean by 7.9\,pp at budget 100 (2 of 3 seeds) and 7.1\,pp at 200 (3 of 3 seeds); a full-budget rerun reproduces the direction (12.3/7.8\,pp), and with warm-started in-domain scorers random still led by 4.2--5.0\,pp at every budget above the shared initial set ($b\geq 50$). A diagnostic explains the mechanism: the synthetic-trained checkpoint scored the public-real pool with fully saturated softmax (mean top1$-$top2 logit margin 100.9 vs.\ $\approx$10.8 in-domain), reducing uncertainty signals to numerical zero---the classical preconditions of uncertainty sampling (a strong scorer, in-domain calibration) are exactly what a low-resource pipeline lacks at startup. This motivates a progressive annotation narrative (annotate first, calibrate the scorer, then select), left to future work. Figure~2 of the supplementary material shows both protocols.

\paragraph{Sensitivity.}
Seed-ratio sweeps show prototype quality saturating early ($\leq$1\,pp beyond 25\% seeds at $K\in\{5,10\}$; $-2.4$\,pp at the best sweep cell $K{=}14$); the confidence threshold $\tau$ and cluster count $K$ interact with the frequency-aware margin design, and full sweeps are registered as future work. The class-imbalance handling switch (frequency-aware margins vs.\ resampling) is registered for evaluation on the public-real tier (not yet run) together with the per-class breakdown mandated by Section~\ref{sec:conclusion}.
